\DocumentMetadata{
  pdfversion=2.0,
  pdfstandard=ua-2,
  lang=en-US,
  tagging=on,
  tagging-setup={math/setup=mathml-SE,tabsorder=structure}
}
\documentclass[11pt,letterpaper]{article}
\usepackage[margin=0.68in,includefoot]{geometry}
\usepackage{newcomputermodern}
\usepackage{amsmath,amscd,mathtools}
\usepackage{tikz,adjustbox}
\providecommand{\square}{\mdlgwhtsquare}
\usepackage[hidelinks]{hyperref}
\hypersetup{
  pdftitle={Analysis PhD Exam, May 1988},
  pdfauthor={University of Florida Department of Mathematics},
  pdfsubject={Graduate mathematics examination},
  pdfdisplaydoctitle=true,
  bookmarksopen=true
}
\setcounter{secnumdepth}{0}
\setlength{\parindent}{0pt}
\setlength{\parskip}{0.28em}
\setlength{\emergencystretch}{3em}
\setlength{\leftmargini}{2.15em}
\setlength{\leftmarginii}{2.65em}
\setlength{\leftmarginiii}{3em}
\pagestyle{empty}
\raggedbottom
\allowdisplaybreaks
\definecolor{ProblemConcern}{HTML}{7A1F1F}
\newenvironment{problemconcern}{%
  \par\tagstructbegin{tag=Aside}\begingroup\color{ProblemConcern}
}{%
  \par\endgroup\tagstructend\nopagebreak[4]\vspace{0.12em}
}
\NewTaggingSocketPlug{block/list/label}{examproblem}{%
  \tagstructbegin{tag=Lbl}\tagstructend%
  \tagstructbegin{tag=\LBody}%
  \tagstructbegin{tag=\ProblemHeadingTag,title-o={\ProblemHeadingText}}%
  \tagmcbegin{tag=\ProblemHeadingTag,actualtext={\ProblemHeadingText}}%
  #2%
  \tagmcend%
  \tagstructend%
}
\newenvironment{examproblems}[2]{%
  \def\ProblemHeadingTag{#2}%
  \AssignTaggingSocketPlug{block/list/label}{examproblem}%
  \begin{enumerate}%
  \setcounter{enumi}{#1}%
  \renewcommand{\labelenumi}{\textbf{\theenumi.}}%
  \setlength{\topsep}{0.35em}%
  \setlength{\partopsep}{0pt}%
  \setlength{\itemsep}{0.48em}%
  \setlength{\parsep}{0.18em}%
}{\end{enumerate}}
\newenvironment{examsubparts}{%
  \AssignTaggingSocketPlug{block/list/label}{default}%
  \begin{enumerate}%
  \setlength{\topsep}{0.18em}%
  \setlength{\partopsep}{0pt}%
  \setlength{\itemsep}{0.16em}%
  \setlength{\parsep}{0.1em}%
}{\end{enumerate}}
\newenvironment{examinstructions}{%
  \par\begingroup\itshape
}{%
  \par\endgroup\vspace{0.18em}
}
\makeatletter
\renewcommand\subsection{\@startsection{subsection}{2}{0pt}%
  {-1.25ex plus -.25ex minus -.1ex}%
  {0.55ex plus .1ex}%
  {\normalfont\large\bfseries}}
\renewcommand{\@maketitle}{%
  \newpage\null\vskip -1em%
  {\footnotesize\bfseries
    \href{https://www.ufl.edu/}{University of Florida}, \href{https://math.ufl.edu/}{Department of Mathematics}\par}%
  \vskip -0.5em%
  \tagstructbegin{tag=H1,title={Analysis PhD Exam, May 1988}}%
  \tagpdfparaOff%
  \tagmcbegin{tag=H1}%
    {\LARGE\bfseries\href{https://gma.math.ufl.edu/exams/analysis-phd/}{Analysis PhD Exam}, May 1988\par}%
  \tagmcend%
  \tagpdfparaOn%
  \tagstructend%
  \vskip 0.25em%
  \tagmcbegin{artifact}%
    \hrule height 1pt%
  \tagmcend%
  \vskip 1em%
}
\makeatother
\title{Analysis PhD Exam, May 1988}
\author{}
\date{}
\begin{document}
\maketitle
\thispagestyle{empty}
\begin{examinstructions}
Answer seven out of the eight questions. Justify all work. To obtain any partial credit, all work must be presented in a neat and logical fashion.
\end{examinstructions}
\begin{examproblems}{0}{H2}
\def\ProblemHeadingText{Problem 1.}
\item
Prove that \(L^1(S,\mathcal A,\mu)\) is a complete space.
\def\ProblemHeadingText{Problem 2.}
\item
State and prove the Hahn–Banach theorem.
\def\ProblemHeadingText{Problem 3.}
\item
State and prove Fubini’s theorem.
\def\ProblemHeadingText{Problem 4.}
\item
Prove that there exists no sequence \((c_n)_{n=1}^{\infty}\) of real numbers such that an infinite series \(\sum a_n\) of real numbers converges if and only if \((c_n a_n)_{n=1}^{\infty}\) is a bounded sequence.

\par
[Hint: Without loss of generality, assume such a sequence exists, with \(c_n\ne 0\) for all~\(n\). Consider the mapping \(T:\ell_\infty\to\ell_1\) given by \(T((x_n))=\frac{x_n}{c_n}\) and use the open mapping theorem.]
\def\ProblemHeadingText{Problem 5.}
\item
Prove that there exists no Banach space of algebraic dimension \(\aleph_0\).

\par
[Hint: show that a finite dimensional subspace is closed and then use the Baire category theorem.]
\def\ProblemHeadingText{Problem 6.}
\item
Let \(f:\mathbb R\to\mathbb R\) be a Lebesgue measurable function. Show that there exists a Borel measurable function~\(g\) such that \(f=g\) a.e. (a.e. with respect to Lebesgue measure).
\def\ProblemHeadingText{Problem 7.}
\item
Let~\(\mu\) be a measure on \(([0,\infty),\mathcal B([0,\infty)))\) defined by
\[
\mu(A)=\int_A x\,dx.
\]
 Let \(T:[0,\infty)\to[0,\infty)\) be given by \(T(x)=e^x-1\). Define a new measure~\(\pi\) on \(([0,\infty),\mathcal B([0,\infty)))\) by setting \(\pi(A)=\mu(T^{-1}(A))\). Compute \(d\pi/d\mu\).
\def\ProblemHeadingText{Problem 8.}
\item
\begin{problemconcern}
\textbf{Warning: Suspected error.} The source visibly defines \(\mu(A)=\int_A x\,dx\) while calling~\(\mu\) a measure on the sigma-algebra of symmetric Borel sets. For sets such as \(\mathbb R\), this signed Lebesgue integral is not defined, so the intended integrand may contain an unprinted or unclear modification. No correction is uniquely determined, and the source reading has been retained.
\end{problemconcern}
Let~\(\lambda\) be Lebesgue measure on \((\mathbb R,\mathcal B(\mathbb R))\). If \(A\in\mathcal B(\mathbb R)\), define the set \(-A=\{x:-x\in A\}\). Define \(\mathcal C=\{A\in\mathcal B(\mathbb R):A=-A\}\).
\begin{examsubparts}
\item[\textnormal{Part (a)}]
Show \(\mathcal C\) is a sigma-algebra.
\item[\textnormal{Part (b)}]
Define two new measures~\(\mu\) and~\(\nu\) on the measurable space \((\mathbb R,\mathcal C)\) by setting
\[
\mu(A)=\int_A x\,dx \qquad\text{whenever }A\in\mathcal C,
\]

\[
\nu(A)=\int_{A\cap[0,\infty)}x\,dx \qquad\text{whenever }A\in\mathcal C.
\]
 Since~\(\lambda\) may also be considered as a measure on \((\mathbb R,\mathcal C)\), we can find two \(\mathcal C\)-measurable functions \(d\mu/d\lambda\) and \(d\nu/d\lambda\) so that
\[
\mu(A)=\int_A \frac{d\mu}{d\lambda}\,dx \qquad\text{whenever }A\in\mathcal C,
\]

\[
\nu(A)=\int_A \frac{d\nu}{d\lambda}\,dx \qquad\text{whenever }A\in\mathcal C.
\]
 Compute these two functions.
\end{examsubparts}
\end{examproblems}
\end{document}
